\section{Supplementary projective calculations}
\label{p1logcbfalt:section}

The reduced-boundary construction of Section~\ref{cyc:section} admits
complementary projective proofs and calculations. We first give
divisorial proofs of the section comparisons and recover the
fixed-source descent by extracting valuations on the base. The fixed
logarithmic zero divisor then gives another discrepancy-sheaf test,
and a toroidal calculation proves logarithmic extension of the root
form. The final subsection tracks the auxiliary subpair when the
minimum-rule divisor is recomputed on a higher base model. This last
operation is distinct from the prescribed boundary modification of
Lemma~\ref{lem:normalized-models}: its discriminant is recomputed,
rather than defined as the strict boundary plus the reduced exceptional
divisor. Throughout these calculations, modifications of the original
source carry strict boundary plus reduced exceptional boundary.

\subsection{Divisorial sections and relative Iitaka systems}

The following elementary comparisons fix the complete systems and the
original models used in the normalization. All varieties in this
subsection are projective over $\C$, all rational divisor degrees are
sufficiently divisible, and the very general fibers are taken integral.
The logarithmic modification identity, also with a pulled-back rational
twist, is \Cref{setup:log-modification}.

\begin{lemma}[Divisorial section test]\label{p1logcbfalt:section-test}
Let \(V\) be normal and \(D\) an integral Weil divisor. A rational section
belongs to \(H^0(V,D)\) if and only if it satisfies the corresponding order
inequality at every prime divisor of \(V\). If \(\pi:V'\to V\) is finite
surjective with \(V'\) normal and \(D\) is Cartier, this test on \(V'\) needs
only primes lying over prime divisors of \(V\). If \(\pi\) is Galois, then
\[
H^0(V',\pi^*D)^{\operatorname{Gal}(V'/V)}=H^0(V,D).
\]
\end{lemma}

\begin{proof}
The first assertion is the description of a rank-one reflexive sheaf as
the intersection of its lattices at codimension-one points. For a finite
morphism between normal varieties, a prime divisor upstairs lies over a
prime divisor downstairs. The last assertion follows from
\((\pi_*\cO_{V'})^{\operatorname{Gal}(V'/V)}=\cO_V\) and the projection
formula.
\end{proof}

\begin{lemma}[Finite pullback]\label{p1logcbfalt:finite}
If \(\pi:V'\to V\) is a dominant generically finite morphism of normal
projective varieties and \(D\) is a rational Cartier divisor on \(V\),
then
\[
\kappa(V',\pi^*D)=\kappa(V,D).
\]
If \(\kappa(V,D)=0\), every nonzero space \(H^0(V,mD)\) has dimension
one.
\end{lemma}

\begin{proof}
Stein factorization reduces the first assertion to a finite morphism,
since a proper birational morphism onto a normal variety has direct image
\(\cO_V\). After replacing \(D\) by a multiple, use the natural inclusion
of graded section rings
\[
R(V,D)\ \subseteq\ R(V',\pi^*D).
\]
Each homogeneous section upstairs is integral over the ring downstairs.
Indeed its monic characteristic polynomial in the finite function-field
extension has coefficients which are rational sections of the appropriate
multiples of \(D\); they are regular in the required lattices at every
prime divisor, and hence are global sections by
Lemma~\ref{p1logcbfalt:section-test}. Thus the two nonzero graded rings have the
same transcendence degree, which gives equality of Iitaka dimensions.
If the ring downstairs has no positive-degree section, the same norm
argument excludes one upstairs. Finally, two independent sections in one
degree have a nonconstant ratio and define a positive-dimensional
rational image. This proves the last assertion.
\end{proof}

\begin{lemma}[Logarithmic pullback]\label{p1logcbfalt:log-pullback}
Let \(\rho:V\to Y\) be a dominant generically finite morphism of smooth projective
varieties. Let \(D_Y\) be a reduced snc divisor, and let \(B\) be a
boundary on \(V\) satisfying \(B\ge(\rho^*D_Y)_{\red}\). Assume that the
boundary supports on the chosen models are snc. Then pullback of rational
pluricanonical forms gives
\[
H^0(Y,m(K_Y+D_Y))\ \longrightarrow\
H^0(V,m(K_V+B))
\]
injectively, for divisible \(m\).
\end{lemma}

\begin{proof}
A logarithmic differential on a smooth snc pair pulls back to a
logarithmic differential along the reduced inverse image. In local
coordinates this follows from
\(d(\rho^*z)/\rho^*z\), whose divisorial poles are simple. Taking top
exterior powers and then tensor powers proves the order inequalities.
Equivalently, at a prime of ramification index \(e\) over a boundary
component, the Jacobian contributes \(e-1\), and the coefficient-one
boundary supplies the remaining one. At exceptional-image primes the
same logarithmic differential calculation applies. Injectivity follows
from dominance and separability.
\end{proof}

\begin{lemma}[Relative Iitaka fibration]\label{p1logcbfalt:iitaka}
Let \(p:V\to B\) be a dominant projective morphism, with \(V\) smooth
projective, and let \(D\) be a rational divisor whose restriction to the
geometric generic fiber has Iitaka dimension \(k\ge0\). After birational
modifications there is a relative Iitaka fibration
\[
V'\xrightarrow{q}U\longrightarrow B
\]
in which \(q\) has geometrically connected generic fiber, such that
\[
\dim U-\dim B=k,\qquad
\kappa(G,D'|_G)=0
\]
for a very general fiber \(G\) of \(q\). Here \(D'=\mu^*D\), where
\(\mu:V'\to V\) is the modification. When \(D=K_V+B_V\), one may
instead take the full log divisor from Lemma~\ref{setup:log-modification}.

If \(D'=K_{V'}+B_{V'}\) and the models are log smooth, then
\(D'|_G\sim_{\Q}K_G+B_G\).
\end{lemma}

\begin{proof}
Apply the Iitaka fibration theorem over the function field of \(B\)
\cite[Section~6, Remark~1]{IitakaOriginal},
take the relative algebraic closure of its section field in \(\C(V)\),
and spread out. Resolve the resulting rational maps and take their
Stein factorizations. The dimension of the relative section field is
\(k\); the restriction of the full divisor to the Iitaka fiber has
Iitaka dimension zero by the Iitaka fibration theorem. For the logarithmic
replacement, the proof of Lemma~\ref{setup:log-modification} gives the sheaf identity
\[
 \mu_*\cO_{V'}\bigl(m(K_{V'}+B_{V'})\bigr)
 =\cO_V\bigl(m(K_V+B_V)\bigr)
\]
in every sufficiently divisible degree. Pushing to the original base
identifies the relative section systems, even if there are no global
sections. Apply the same full-divisor Iitaka theorem to the replacement
divisor; its relative section field identifies its fibration with the
one already constructed and gives the asserted Iitaka dimension zero
on the fiber.
Generic smoothness and adjunction give the last assertion.
The dimensions at very general complex fibers agree with the geometric
generic values: for each divisible degree use cohomology and base
change, and remove the union of the exceptional closed sets over the
countably many degrees.
\end{proof}

\begin{lemma}[Easy addition]\label{p1logcbfalt:easy-addition}
Let \(p:V\to B\) be dominant and projective, with \(V\) smooth projective
and integral very general fiber \(G\), and let \(D\) be a rational
divisor on \(V\). Then
\[
\kappa(V,D)\le \dim B+\kappa(G,D|_G).
\]
If \(D\) has a nonzero global section, its restriction to a general
fiber is nonzero.
\end{lemma}

\begin{proof}
A nonzero section cannot vanish identically on all fibers over a dense
open subset. For each divisible \(m\) with sections, let \(\phi_m\)
be its rational map. The image of \((p,\phi_m)\) has dimension at most
\(\dim B+\kappa(G,D|_G)\), and dominates the image of \(\phi_m\).
Taking the maximum over \(m\) proves the inequality. If the fiber term
is \(-\infty\), the last assertion implies that the total term is
\(-\infty\) as well.
\end{proof}

\subsection{Restricted valuations and the fixed original source}

Let $p:H\to I$ be a dominant morphism of smooth projective varieties,
and fix the model $H$. A divisorial valuation $v$ of $\C(H)$ whose
restriction to $\C(I)$ is nontrivial restricts to a positive integral
multiple of a divisorial valuation of $\C(I)$. To see this, put
$n=\dim H$, $m=\dim I$, and denote the restricted valuation by $w$.
The value groups have rational rank one with finite-index inclusion.
The relative residue-transcendence-degree inequality gives
\[
 \operatorname{trdeg}_{\kappa(w)}\kappa(v)\le n-m.
\]
As $\operatorname{trdeg}_{\C}\kappa(v)=n-1$, this implies
$\operatorname{trdeg}_{\C}\kappa(w)\ge m-1$. The Abhyankar
inequality gives the reverse inequality. Thus $w$ is divisorial.
Finitely many such restrictions can be extracted on one smooth
projective birational model of $I$.

Only finitely many prime divisors of the fixed $H$ have image of
codimension at least two in $I$: they lie in the complement of a
locus where $p$ is smooth and are components of its divisorial part.
Extract their restrictions as above. Original horizontal primes
remain horizontal, while original vertical primes then have
divisorial centers on the prepared base. These centers remain
divisorial on higher base models. Hence on any subsequent commuting
birational diagram, a source prime whose base image still has
codimension at least two is exceptional over the fixed $H$.

For the rank-one section comparison, fix the original smooth pair
$(X,D_X)$, and let $X'\to X$ carry strict plus reduced exceptional
boundary $D'$. Let $x:X'\to W$ be the relative Iitaka fibration
on smooth projective models; its geometric generic fiber $(J,D_J)$
has logarithmic Kodaira dimension zero. Let $s$ be a nonzero
relative rational $p$-canonical form generating the degree-$p$
generic section space. For a source prime $E$, put
$d_E=\operatorname{coeff}_E D'$, and for a base prime $P$ define
\[
 t_P=\min_{E\mapsto P}
 \frac{p^{-1}\operatorname{ord}^{pK_{X'/W}}_E(s)+d_E}
 {\operatorname{ord}_E(x^*P)},\qquad
 T_W=\sum_Pt_PP.
\]
Here $T_W$ is the entire minimum-rule base term, before its division
into discriminant and moduli parts. In divisible degrees, division by
$s^{m/p}$ gives
\begin{equation}\label{p1logcbfalt:source-descent}
 H^0\bigl(X',m(K_{X'}+D')\bigr)
 \lhook\joinrel\longrightarrow
 H^0\bigl(W,m(K_W+T_W)\bigr).
\end{equation}
Indeed, on the geometric generic fiber the section is a constant
multiple of $s^{m/p}$; connectedness identifies that coefficient
with a rational $m$-canonical form $\sigma$ on $W$. At a component
$E$ dominating $P$, write $\sigma$ in a nonvanishing local frame of
$mK_W$, and put $a_E=\operatorname{ord}_E(x^*P)$. Its logarithmic
regularity upstairs is equivalent to
\[
 a_E\operatorname{ord}_P(\sigma)
 +(m/p)\operatorname{ord}^{pK_{X'/W}}_E(s)+md_E\ge0.
\]
The inequalities for all components above $P$ are exactly
$\operatorname{ord}_P(\sigma)+mt_P\ge0$. Horizontal primes impose
no further condition because $s$ is an allowed generic logarithmic
pluriform. The same argument works over the generic point of the
original base $Y$. If the relative Iitaka dimension is $k$, then
$\dim(W/Y)=k$, so preservation of section ratios makes
$K_W+T_W$ big over $Y$.

For the converse, extract on $W$ the restrictions of all source
divisors with exceptional image that survive as divisors on the
original $X$. Resolve the induced diagram and recompute the minimum
rule. Every divisor of the original $X$ is now tested either by the
generic horizontal condition or by a base prime. A section on the
right of \eqref{p1logcbfalt:source-descent}, multiplied by $s^{m/p}$,
therefore satisfies the logarithmic order inequality at every
original source prime. Any remaining untested pole lies on a
divisor exceptional over $X$. The codimension-one criterion on the
normal original source gives a section of $m(K_X+D_X)$; later
models preserve this criterion. Conversely the logarithmic
modification identity recovers its section on $X'$. The
identifications respect powers and multiplication. Under these
identifications, corresponding section ratios agree after pullback by
$x:X'\to W$. Their section fields and rational-map image dimensions
therefore agree.

\subsection{The fixed zero divisor and the discrepancy-sheaf test}

In the notation of Lemma~\ref{lem:reduced-cyclic-top}, the rational
divisor $Z=E/p$ is independent of the chosen nonzero logarithmic
pluriform and its degree. Indeed, if $\omega_m$ is a nonzero section
in another degree $m$, the degree-$mp$ sections $\omega_m^p$ and
$\omega^m$ are proportional, since that section space is
one-dimensional. Consequently
\[
 p\bigl(\operatorname{div}(\omega_m)+mD\bigr)=mE.
\]
The least positive integer trivializing $K_J+D-Z$ is also $p$.
Writing $\omega=b\eta^{\otimes p}$ in a rational canonical frame
gives $p(K_J+D-Z)=-\operatorname{div}(b)$. Conversely, a
trivialization of $r(K_J+D-Z)$ gives a rational $r$-canonical form
whose logarithmic zero divisor is $rZ\ge0$, hence a nonzero
logarithmic pluriform of degree $r$. Minimality gives $r\ge p$.

Here is a second way to check the generic discrepancy-sheaf rank
condition in Corollary~\ref{cor:reduced-integral-moduli}. Take a log
resolution $J'\to J$, let $D_{J'}$ be the strict boundary plus the
reduced exceptional divisor, and put
\[
 E_{J'}=\operatorname{div}(\omega|_{J'})+pD_{J'}\ge0.
\]
The crepant transform of the auxiliary generic subboundary is
$D_{J'}-E_{J'}/p$. Its coefficients are at most one. If $A^*$ is
the discrepancy b-divisor with coefficient $-1$ replaced by zero,
the rounded trace is effective, and a positive coefficient can
occur only on $\operatorname{Supp}E_{J'}$. Coefficientwise,
\[
 0\le \lceil A^*\rceil_{J'}
   \le\lceil E_{J'}/p\rceil\le E_{J'}.
\]
The first inequality uses the omission of discrepancy $-1$.
Logarithmic birational invariance gives $\kappa(E_{J'})=0$.
The associated space of rational functions therefore has dimension
at most one and contains the constants. This proves rank one on
every sufficiently high resolution, including when the auxiliary
subboundary has negative coefficients.

\subsection{A toroidal proof of logarithmic extension}

The extension part of Lemma~\ref{lem:reduced-cyclic-top} also has a
direct coordinate proof on a toroidal resolution of the SNC
cyclic-cover construction. First resolve the supports on $J$,
retaining strict plus reduced exceptional logarithmic boundary.
The minimal index and the fixed-zero assertions are unchanged.
In a local SNC chart write
\[
 \omega=\varepsilon\prod_i z_i^{a_i-pd_i}
                (dz_1\wedge\cdots\wedge dz_r)^{\otimes p},
 \qquad a_i\ge0,\quad d_i\in\{0,1\},
\]
where $\varepsilon$ is a unit. Include all coordinates; outside
$D+\operatorname{Supp}E$ put $d_i=a_i=0$. After a local unramified
change the unit has a root and contributes no divisorial order.
Let $v$ be a toroidal prime on the resolved root cover, and put
$n_i=v(z_i)\ge0$. The nonzero ray defining this prime has at least
one positive $n_i$. Rewriting with logarithmic differentials gives
\begin{equation}\label{p1logcbfalt:toroidal-root-order}
 \operatorname{ord}_v(\Omega)
   =\sum_i n_i\left(1-d_i+\frac{a_i}{p}\right)-1.
\end{equation}
The logarithmic volume form has order $-1$ along a toroidal prime,
and the change of torus lattice multiplies it by a nonzero constant.
Every coefficient in the sum is nonnegative. It vanishes exactly
when $d_i=1$ and $a_i=0$, namely along the coefficient-one pole
boundary of the auxiliary subpair.

If $v$ is outside the inverse image of this pole boundary, then
$n_i=0$ for each such zero coefficient. Some remaining $n_i$ is
positive, so the sum in \eqref{p1logcbfalt:toroidal-root-order} is
strictly positive. The order of the rational top form $\Omega$ is
integral, so that positive sum is at least one, giving regularity.
On the inverse pole boundary the nonnegative sum gives order at
least $-1$. Primes in the smooth unramified locus have the
regularity already checked on the finite normalization. Thus
$\Omega$ extends as a top form with logarithmic poles allowed only
along the inverse pole boundary. This provides the coordinate
alternative to the discrepancy-scaling proof.

\subsection{The normalized discriminant on higher models}

Let $x:X'\to W$ and $D'$ be the projective reduced-boundary
data of Lemma~\ref{lem:rankone-normalization}. Write $s$ for its
relative rational $p$-canonical generator, and set
\[
 t_P=\min_{E\mapsto P}
 \frac{p^{-1}\operatorname{ord}^{pK_{X'/W}}_E(s)+d_E}
 {\operatorname{ord}_E(x^*P)},\qquad
 T_W=\sum_Pt_PP,
\]
\[
 \Delta_0=-p^{-1}\operatorname{div}_{K_{X'/W}}(s),
 \qquad \Delta=\Delta_0+x^*T_W.
\]
In this subsection $T_W$ denotes the entire minimum-rule base
term, before its decomposition into discriminant and moduli parts;
the boundary denoted $T$ in Section~\ref{rankone:section} is the
discriminant part. We have
$K_{X'}+\Delta\sim_{\Q}x^*(K_W+T_W)$.
The auxiliary divisor is required to be sub-lc over the generic
point of $W$. Its coefficients at source divisors with image of
codimension at least two are not asserted to be at most one.

For a prime $P\subset W$, let
\[
 c_P=\operatorname{lct}_{\eta_P}(X',\Delta;x^*P),
 \qquad b_P=1-c_P,
 \qquad B_W=\sum_Pb_PP.
\]
Over the generic point of $P$ the minimum rule gives
$\Delta\le D'$, including the horizontal divisors. Log canonicity
therefore gives $c_P\ge0$. A component $E$ attaining the minimum
has coefficient exactly $d_E\in\{0,1\}$ in $\Delta$, and hence
\[
 c_P\le\frac{1-d_E}{\operatorname{ord}_E(x^*P)}\le1.
\]
It follows that $0\le B_W\le1$. If $P$ lies in a required inverse
boundary from a base $Y$, all components over $P$ have $d_E=1$,
so $c_P=0$ and $b_P=1$.

For the exact model comparison, let $\beta:W_1\to W_0$ be a
higher smooth base model and let $\pi:X_1\to X_0$ resolve the
induced diagram, with $x_i:X_i\to W_i$. Choose compatible
rational top forms for the canonical divisors. First form the
crepant transform
\[
 K_{X_1}+\Delta_{\mathrm{crepant}}
    =\pi^*(K_{X_0}+\Delta_{\mathrm{old}}).
\]
Using the new logarithmic boundary and the same relative
trivialization $s$, recompute the minimum-rule divisor $T_{W_1}$
and the normalized auxiliary divisor $\Delta_{\mathrm{new}}$.
Comparison of their definitions gives
\begin{equation}\label{p1logcbfalt:model-twist}
 \begin{split}
 Q&=K_{W_1}+T_{W_1}-\beta^*(K_{W_0}+T_{W_0}),\\
 \Delta_{\mathrm{new}}
   &=\Delta_{\mathrm{crepant}}+x_1^*Q.
 \end{split}
\end{equation}
Thus the twist includes the canonical discrepancy of the base.
Adding $x^*Q$ lowers the threshold at $P$ by
$\operatorname{coeff}_P Q$, adding $Q$ both to the discriminant
and to the total base term. It leaves the moduli part unchanged.
Birational pullback of the lc-trivial fibration also preserves its
moduli b-divisor; these compatibilities agree with the Hodge
extension description in \cite[Remarks~6.22 and~6.25]{BFMT}.

Resolve the finite supports on $W$ and resolve the induced source
diagram. At an unchanged codimension-one point, logarithmic forms
satisfying the old tests remain logarithmic on the source
resolution, while the strict old minimizing components persist.
Thus the minimum and discriminant trace there are unchanged.
New support is contained in the base-exceptional divisor. Its union
with the old strict support is SNC on the chosen resolution;
recomputing and applying the preceding coefficient argument yields
an effective SNC discriminant with the same coefficient-one inverse
boundary and the same moduli b-divisor.

Finally, if $m=ap$ and $s_m=c s^a$ on the generic fiber, with
$c\in\C(W)^*$, the minimum construction gives
\[
 T_W(s_m)=T_W(s)+m^{-1}\operatorname{div}(c).
\]
The two occurrences of $\operatorname{div}(c)$ cancel in the
definition of the normalized auxiliary subpair. It is unchanged.
The minimal-index cover is therefore compatible with all common
divisible degrees used for section descent.
